% !TEX root = ../../main.tex
\section{Covariance matrix and correlation of a joint PMF}

% (a) Marginal PMFs
% (b) Conditional PMF
% (c) Mean vector
% (d) Covariance matrix
% (e) Correlation coefficient

The PMF of the random vector $\vec{X} = [X_1, X_2]$ is given in the \autoref{tab:joint-pmf} which is visualized in \autoref{fig:joint-pmf}. 
\begin{table}[h]
    \centering
    \caption{Joint PMF $p_{\vec{X}}[x_1, x_2]$}\label{tab:joint-pmf}
    \begin{tabular}{c|ccc}
        & $x_2 = 0$ & $x_2=1$ & $x_2=2$ \\ \hline
        $x_1=0$ & $\frac{1}{8}$ & $0$ & $\frac{1}{4}$ \\
        $x_1=1$ & $0$ & $\frac{1}{8}$ & $\frac{1}{4}$ \\
        $x_1=2$ & $\frac{1}{8}$ & $0$ & $\frac{1}{8}$
    \end{tabular}
\end{table}
The Sample Space of the Random Variables is
\[
    S_{X_1} = \{0,1,2\} \qquad S_{X_2} = \{0,1,2\} \qquad S_{\vec{Y}} = S_{X_1} \times S_{X_2}
\]
\begin{figure}[h]
    \centering
    \includegraphics[width=0.4\textwidth]{JointPMF_L.pdf}
    \caption{Joint PMF $p_{\vec{X}}[x_1, x_2]$}\label{fig:joint-pmf}
\end{figure}

\subsection{Marginal PMFs}

The Marginal PMFs $p_{X_1}[x_{1}]$ and $p_{X_2}[x_{2}]$ are calculated with \autoref{eq:marginal-px1} and \autoref{eq:marginal-px2} respectively.
\begin{align}
    p_{X_1}[x_{1}] = \sum_{x_{2} \in S_{X_2}} p_{\vec{X}}[x_1, x_2] \label{eq:marginal-px1}\\
    p_{X_2}[x_{2}] = \sum_{x_{1} \in S_{X_1}} p_{\vec{X}}[x_1, x_2] \label{eq:marginal-px2}
\end{align}

Meaning for each instance of $x_n$, we sum over all the other probabilities. This can be best show by reconsidering the PMF and summing over the rows or columns in \autoref{tab:marginal-pmf}.

\begin{table}[h]
    \centering
    \caption{Calculation of Marginals}\label{tab:marginal-pmf}
    \begin{tabular}{c|ccc|c}
        & $x_2 = 0$ & $x_2=1$ & $x_2=2$ & $p_{X_1}[x_1]$ \\ \hline
        $x_1=0$ & $\tfrac{1}{8}$ & $0$ & $\tfrac{1}{4}$ & $\tfrac{3}{8}$ \\
        $x_1=1$ & $0$ & $\frac{1}{8}$ & $\tfrac{1}{4}$ & $\tfrac{3}{8}$ \\
        $x_1=2$ & $\tfrac{1}{8}$ & $0$ & $\tfrac{1}{8}$ & $\tfrac{1}{4}$ \\ \hline
        $p_{X_2}[x_2]$ & $\tfrac{1}{4}$ & $\tfrac{1}{8}$ & $\tfrac{5}{8}$
    \end{tabular}
\end{table}

The marginals are plotted in \autoref{fig:marginals}.

\begin{figure}[h]
    \centering
    \begin{subfigure}[c]{0.3\textwidth}
        \centering
        \includegraphics[width=\linewidth]{MarginalPMF-X2.pdf}
        \caption{$p_{X2}$}
    \end{subfigure}\hfill$\impliedby$\hfill
    \begin{subfigure}[c]{0.3\textwidth}
        \centering
        \includegraphics[width=\linewidth]{JointPMF.pdf}
        \caption{Joint}
    \end{subfigure}\hfill$\implies$\hfill
    \begin{subfigure}[c]{0.3\textwidth}
        \centering
        \includegraphics[width=\linewidth]{MarginalPMF-X1.pdf}
        \caption{$p_{X1}$}
    \end{subfigure}
    \caption{Marginal Distributions}\label{fig:marginals}
\end{figure}

\subsection{Conditional PMF}

The Conditional PMF $p_{X_2|X_1}[x_{2}|x_{1}]$ is given by \autoref{eq:conditional}.
\begin{equation}
    p_{X_2|X_1}[x_{2}|x_{1}] = \frac{p_{\vec{X}}[x_1, x_2]}{p_{X_1}[x_1]} \label{eq:conditional}
\end{equation}

From \autoref{tab:marginal-pmf} we know, that $p_{X_1}[x_1 = 1] = \frac{3}{4}$. Finally, we take the row in the PMF where $x_1=1$ and get
\[
    p_{X_2|X_1}[x_{2}|x_{1}=1] = \frac{p_{\vec{X}}[x_1=1, x_2]}{p_{X_1}[x_1=1]} = \begin{bmatrix}
        0 & \tfrac{1}{3} & \tfrac{2}{3}
    \end{bmatrix}
\]

\begin{figure}
\centering
\hfill
\begin{minipage}{.4\textwidth}
    \centering
    \includegraphics[width=0.9\textwidth]{ConditionalPMF-X2X1.pdf}
    \caption{Conditional PMF, where $x_1$ is given as $x_1 = 1$.}\label{fig:conditional-pmf}
\end{minipage}\hfill
\begin{minipage}{.4\textwidth}
    \centering
    \includegraphics[width=0.9\textwidth]{PMFWithMeanStem.pdf}
    \caption{Joint PMF with the mean vector plotted}\label{fig:pmf-with-mean}
\end{minipage}
\hfill
\end{figure}

\subsection{Mean Vector}

The mean is the first statistical order, which for each random variable in the random vector $\vec X$ is given by:
\[
\mu_{X_n} = \E{X_n} \DefinedBy \sum_{x_n \in S_{X_n}} x_n p_{X_n}[x_n]
\]

The complete mean vector is then 
\[
\vec{\mu}_{\vec{X}} = \begin{bmatrix}
    \mu_{X_1} \\
    \mu_{X_2}
\end{bmatrix} = \begin{bmatrix}
    \tfrac{7}{8} \\
    \tfrac{11}{8}
\end{bmatrix}
\]

\subsection{Auto-Covariance Matrix}

The Auto-Covariance Matrix is defined by
\begin{align}
    \mat{C}_{\vec{XX}}
    &= \E{(\vec{X} - \E{X}){(\vec{X} - \E{X})}^{\T}} \notag \\
    &= \E{(\vec{X} - \vec{\mu}_{\vec{X}}){(\vec{X} - \vec{\mu}_{\vec{X}})}^{\T}} \label{eq:auto-cov-mtx}
\end{align}
For easier Calculations, we normalize $\vec{X}$ in the first order moment.
\[
\vec{Y} = \vec{X} - \vec{\mu}_{\vec{X}} = \begin{bmatrix}
    X_1 - \mu_{X_1} \\
    X_2 - \mu_{X_2}
\end{bmatrix}
\]
so the Sample space and the PMF (see \autoref{tab:joint-pmf-y}) become
\[
    S_{Y_1} = \{-\tfrac{7}{8}, \tfrac{1}{8}, \tfrac{9}{8}\} \qquad S_{Y_2} = \{-\tfrac{11}{8}, -\tfrac{3}{8}, \tfrac{5}{8}\} \qquad S_{\vec{Y}} = S_{Y_1} \times S_{Y_2}
\]
\begin{table}[h]
    \centering
    \caption{Joint PMF $p_{\vec{Y}}[y_1, y_2]$}\label{tab:joint-pmf-y}
    \begin{tabular}{c|ccc}
        & $y_2 = -\tfrac{11}{8}$ & $y_2=-\tfrac{3}{8}$ & $y_2=\tfrac{5}{8}$ \\ \hline
        $y_1=-\tfrac{7}{8}$ & $\frac{1}{8}$ & $0$ & $\frac{1}{4}$ \\
        $y_1=\tfrac{1}{8}$ & $0$ & $\frac{1}{8}$ & $\frac{1}{4}$ \\
        $y_1=\tfrac{9}{8}$ & $\frac{1}{8}$ & $0$ & $\frac{1}{8}$
    \end{tabular}
\end{table}

Now, the PMF has its center of mass around the origin, because the mean vector of $\vec{Y}$ is zero. To calculate the Auto-Covariance Matrix, we expand the 
\begin{align*}
    \mat{C}_{\vec{XX}} = \E{\vec{Y}\vec{Y^{\T}}}
    &= \sum_{y_1 \in S_{Y_1}}\sum_{y_2 \in S_{Y_2}} \vec{yy}^\T p_{\vec{Y}}[y_1,y_2] \\
    &= \sum_{y_1 \in S_{Y_1}}\sum_{y_2 \in S_{Y_2}} [y_1 y_2][y_1 y_2]^\T p_{\vec{Y}}[y_1,y_2] \\
    &= \sum_{y_1 \in S_{Y_1}}\sum_{y_2 \in S_{Y_2}} \begin{bmatrix}
        y_1y_1 & y_1y_2 \\ y_2y_1 & y_2y_2
    \end{bmatrix} p_{\vec{Y}}[y_1,y_2] 
\end{align*}
So for each element in the auto-covariance Matrix, we get
\[
    \E{Y_i Y_j} = \sum_{y_1 \in S_{Y_1}}\sum_{y_2 \in S_{Y_2}} y_i y_j p_{\vec{Y}}[y_1,y_2]
\]
Computing the Covariance Matrix in MATLAB we receive
\begin{equation}
    \mat{C}_{\vec{XX}} = \begin{bmatrix}
        0.609375 & -0.078125 \\
        -0.078125 & 0.734375
    \end{bmatrix} = \begin{bmatrix}
        \sfrac{39}{64} & \sfrac{-5}{64} \\
        \sfrac{-5}{64} & \sfrac{47}{64}
    \end{bmatrix}
\end{equation}

\subsection{Correlation Coefficient}

The correlation coefficient is given by 
\begin{equation}
    \rho_{X_1,X_2} = \frac{\Cov{X_1}{X_2}}{\sqrt{\Var{X_1}\Var{X_2}}}
\end{equation}
because 
\begin{equation}
    \mat{C}_{\vec{XX}} = \begin{bmatrix}
        \Var{X_1} & \Cov{X_1}{X_2} \\
        \Cov{X_2}{X_1} & \Var{X_2}
    \end{bmatrix}
\end{equation}
We can use the values from the auto correlation matrix and calculate the correlation coefficient straight away
\[
    \rho_{X_1,X_2} = -\frac{\tfrac{5}{64}}{\sqrt{\tfrac{39}{64}\tfrac{47}{64}}} = -\frac{5}{\sqrt{39\cdot47}} \approx -0.116785
\]